\cotefile{156-8}
\scanpage{82}
\mspage{73}
% archive sheet number: 82

\deleted{\unsure{Obs.|Soit.} Supposons}

% The following opening paragraph is cancelled by three long diagonal strokes.
\deleted{Soit $F\inserted{\subset M}$ une figure de $M$. La chose \unsure{essentielle|naturelle}, ici, c'est que}
\[
\deleted{X,Y\in F\Longrightarrow X\unsure{\gtrless|\asymp}Y}
\]
\deleted{i.e. que $F$ soit \unsure{contenu|contenue} dans une figure, (et non \inserted{\unsure{tellement|seulement}} que \deleted{\unsure{elle}} soit \deleted{elle} une figure vraiment, i.e. fermée pour $\leq$).}
% Both occurrences of "elle" also have local strike-throughs inside the cancelled block.

Je vais m'intéresser d'abord, plus particulièrement, à
\[
\operatorname{supp}^{\circ}X=A\subset M\qquad(X\in M)
\]
et sa relation à
\[
|X|^{\circ}\subset P.
\]
On a $X\in\operatorname{supp}^{\circ}X$ \unsure{donc|d'où}
\[
|X|^{\circ}\subset|\operatorname{supp}^{\circ}X|^{\circ}
\]
j'aimerais que l'\inserted{inclusion} inverse soit vrai, i.e. que
\[
|X|^{\circ}=|\operatorname{supp}^{\circ}X|^{\circ}
\]
ou encore que
\[
\forall Y\in\operatorname{supp}^{\circ}X,\quad\text{on ait }|Y|^{\circ}\subset|X|^{\circ}
\]
(l'implication inverse étant claire).
Il revient au $=$ dire que
\[
\begin{aligned}
&\deleted{Y\notin\unsure{\operatorname{supp}|\operatorname{supp}^{\circ}}}\\[-2pt]
&|Y|^{\circ}\not\subset|X|^{\circ}\Longrightarrow Y\notin\operatorname{supp}^{\circ}X,\quad\text{i.e.}\\
&\text{i.e. }|Y|^{\circ}\cap\complement|X|^{\circ}\neq\varnothing
\qquad\begin{gathered}
\exists Z\in M\ \unsure{\text{avec}|\text{un}}\\
Z\mathrel{|\circ|}X,\quad Z\mathrel{\overline{|\circ|}}Y.\\
\text{i.e.}\\
\left\{\begin{aligned}
|Z|^{\circ}\cap|X|^{\circ}&=\varnothing\\
|Z|^{\circ}\cap|Y|^{\circ}&\neq\varnothing
\end{aligned}\right.
\end{gathered}
\end{aligned}
\]
\clearpage
\scanpage{83}
\mspage{74}
% archive sheet number: 83

C'est clair dans le cas où
\[
|Y|^{\circ}\cap\partial X\neq\varnothing,
\]
\deleted{i.e. quand $\exists Z\leq X$, $\unsure{Z}$ \unsure{avec}}

ce qui signifie, puisque
\[
\partial X=\bigcup_{Z<X}|Z|^{\circ},
\]
que
\[
\exists Z<X,\quad |Y|^{\circ}\cap|Z|^{\circ}\neq\varnothing,
\]
on prend alors ce $Z$-là. \unsure{Sinon|disons}, on aurait donc
\[
|Y|^{\circ}\deleted{\unsure{\cap}}\cap\complement|X|\neq\varnothing
\]
\[
\text{puisque }\complement|X|^{\circ}=|\partial X|\amalg\complement|X|,
\]
et on gagne si on fait l'hypothèse suivante
\[
\begin{aligned}
\text{\uline{Hyp}}\quad&\forall X\in M,\ \deleted{\text{\unsure{on}}}\text{ et tout }x\in P\setminus|X|,\\
&\left|\begin{aligned}
&\exists Z\in M\deleted{\unsure{\complement|X|}}\quad\text{avec}\\
&Z\parallel X,\qquad x\in|Z|
\end{aligned}\right.\\
&\text{(i.e. }|Z|\cap|X|=\varnothing\text{)}
\end{aligned}
\]
\margin{\uline{NB} quitte à remplacer $Z$ par une de ses strates, \unsure{OPS|OQS} $\unsure{\Leftarrow|\Leftrightarrow}\ x\in|Z|^{\circ}$}
% slanted; note connected to the hypothesis by a short arrow
\margin{$\uparrow$ C'est trop fort, il est plus raisonnable d'écrire $Z\mathrel{|\circ|}X$ (pour inclure p.ex. les \unsure{magmiens|magmas} fidèles, et notamment les maquettes)}

Pour ce qui m'occupe à présent, il suffirait d'avoir $\deleted{\unsure{|Z|}}Z\mathrel{|\circ|}X$, i.e. $|Z|^{\circ}\cap|X|^{\circ}=\varnothing$.

Cette hypothèse équivaut \unsure{(pour $\alpha$ donné à)|(pour $X$ donné à)}
\[
\begin{aligned}
P\setminus|X|&=\bigcup_{Z\in\operatorname{cosupp}^{\unsure{\sharp|\parallel}}(X)}|Z|\qquad\deleted{\unsure{\text{ou}}}\\
&\left(=\bigcup_{Z\in\operatorname{cosupp}^{\sharp}(X)}|Z|^{\circ}\right)
\end{aligned}
\]
\clearpage
\scanpage{84}
\mspage{75}
% archive sheet number: 84

\deleted{En fait, pour le moment \unsure{il me suffirait|il me suffit} d'avoir}
Ceci implique
\[
P\setminus|X|^{\circ}=\bigcup_{Z\in\operatorname{cosupp}^{\circ}X}|Z|^{\circ}
\]
\margin{\uline{NB} Avec cette seule relation, il n'est pas exclu que parmi les $Z$ qui interviennent, il y en ait avec $|Z|\cap|X|^{\circ}\neq\varnothing$ !}
et ceci me suffit pour l'instant. On voit que ça assure déjà les deux relations
\[
\begin{aligned}
\unsure{\varphi|\phi}(\operatorname{supp}^{\circ}X)&\overset{\mathrm{def}}{=}\unsure{\lvert}\operatorname{supp}^{\circ}X|^{\circ}=|X|^{\circ}\\
\varphi(\operatorname{cosupp}^{\circ}X)&\overset{\mathrm{def}}{=}|\operatorname{cosupp}^{\circ}X|^{\circ}=P\setminus|X|^{\circ}.
\end{aligned}
\]

\deleted{qui impliquent} Je dis qu'on \unsure{aura|a} aussi
\[
\begin{aligned}
&\deleted{\operatorname{supp}^{\circ}X}\\[-2pt]
\varphi(\operatorname{supp}X)&=|X|\\
\varphi(\operatorname{cosupp}X)&=P\setminus|X|
\end{aligned}
\]
\unsure{où|si} \deleted{\unsure{supp}} par définition
\[
\begin{aligned}
\operatorname{supp}X&=\operatorname{supp}^{\circ}(\widetilde X)\\
\operatorname{cosupp}X&=\operatorname{cosupp}^{\circ}(\widetilde X)
\end{aligned}
\]
Comme \deleted{\unsure{pour}} $\forall Y\in\widetilde X$, on a $Y\in\operatorname{supp}X$, d'où $\deleted{\varphi(Y)\subset}|Y|^{\circ}\subset\varphi(\operatorname{supp}X)$, on voit que $\varphi(\operatorname{supp}X)$ contient les $|Y|^{\circ}$, donc leur réunion $|X|$
\[
|X|\subset\varphi(\operatorname{supp}X)
\]
Pour l'inclusion inverse, il faut voir que
\[
Y\in\operatorname{supp}X\bigl(=\operatorname{supp}^{\circ}(\widetilde X)\bigr)\Longrightarrow|Y|^{\circ}\subset|X|
\]
ou encore, comme tantôt
\[
\begin{aligned}
|Y|^{\circ}\not\subset|\unsure{\overline X|\widetilde X}|\quad&\Longrightarrow\quad Y\notin\operatorname{supp}^{\circ}(\widetilde X)\quad\text{i.e. }\exists Z\\
\text{i.e. }|Y|^{\circ}\cap(P\setminus\overline X)\neq\varnothing\quad&\qquad\text{avec }\widetilde X\mathrel{|\circ|}Z,\quad Y\mathrel{|\circ|}Z,\\
&\text{i.e. }|X|\cap|Z|^{\circ}\neq\varnothing,\quad |Y|^{\circ}\cap|Z|^{\circ},
\end{aligned}
\]
% The last written line reaches the physical lower edge of the supplied image.
\clearpage
\scanpage{85}
\mspage{76}
% archive sheet number: 85

C'est \inserted{donc} ici qu'on aura besoin de l'hypothèse plus haut, \unsure{des|les} \deleted{\unsure{deux}}, \unsure{de fait|à faire}.$^{*}$
% slanted
\margin{$^{*}$ Pas tout à fait puisqu'on a besoin \unsure{pour ceci|par ici}, des $Z$ tels que $|Z|\subset P\setminus|X|$ i.e. $|Z|\cap|X|=\varnothing$, or on a seulement $|Z|^{\circ}\subset P\setminus|X|$.}

Mais je suspecte que ce n'est pas encore tout à fait vrai --- comme exemple, Soit $Y<X$, et considérons
\[
\operatorname{supp}^{\circ}(\{Y,X\})\quad\text{et}\quad\unsure{\varphi|\phi}(\operatorname{supp}^{\circ}(\{Y,X\}))
\]
on aura donc
\[
\deleted{\unsure{\{Y\}}}\quad |X|^{\circ}\cup|Y|^{\circ}\subset\varphi(\operatorname{supp}^{\circ}(\{Y,X\}))
\]
aura-t-on l'inverse ? \deleted{Ici, pour} \inserted{Mais oui, ---} et je vois \inserted{à présent} quelle est l'hypothèse générale qui \deleted{\unsure{assure|entraîne}} \inserted{\unsure{assurera|assure}} qu'on a un ``bon comportement'' pour des ``supports \inserted{(\unsure{intérieurs|internes})} constructibles''.

\uline{Définition} Soit $\unsure{\mathscr F|\mathcal F}$ un \unsure{atlas} associé au magmien $M$. Un élément \inserted{$A$} de $\overline\Sigma_M$ est dit un support (\unsure{interne|intérieur}) \uline{constructible}, s'il existe une \deleted{figure $F\in\mathscr F$ et une} partie $\unsure{\Phi|\mathscr P}$ de $M$, majorée dans $\mathscr F$, \inserted{(i.e. $\Phi\subset F$, avec $F\in\mathscr F$)} telle que l'on ait
\[
A=\deleted{\unsure{\Phi|F}}\operatorname{supp}^{\circ}(\Phi).
\]
On dit qu'une fonction support ensembliste externe [fidèle]
\[
X\longmapsto|X|,\qquad M\longrightarrow\unsure{\mathscr P|\mathcal P}(P)
\]
est \uline{pleinement fidèle} (pour $\mathscr F$), si elle
\clearpage
\scanpage{86}
\mspage{77}
% archive sheet number: 86

\begin{samepage}
satisfait la propriété suivante
\[
\forall F\in\unsure{\mathscr F|\mathcal F},\quad\text{on a}\quad
\unsure{\varphi|\phi}(\operatorname{cosupp}^{\circ}(F))=\tikz[remember picture,baseline=(s86xone.base)]{\node[draw,circle,inner sep=1pt] (s86xone) {$X$};}\setminus|F|
\]
\reversemarginpar
\margin{\tikz[remember picture]{\node[inner sep=0pt] (s86note) {$?$ \uline{plutôt} $P$};}}
\normalmarginpar
% The marginal correction has arrows to the first two circled X's and the X after the quantifier in the boxed line.
i.e.
\[
\begin{aligned}
\tikz[remember picture,baseline=(s86xtwo.base)]{\node[draw,circle,inner sep=1pt] (s86xtwo) {$X$};}\setminus|F|&=\bigcup_{\substack{Z\in M\\\text{t.q. }\tikz[remember picture,baseline=(s86cond.base)]{\node[inner sep=0pt] (s86cond) {$\scriptstyle Z\mathrel{|\circ|}F$};}}}|Z|^{\circ}\\
&\bigl[\text{i.e. }|Z|^{\circ}\subset X\setminus|F|\bigr]
\end{aligned}
\]
\margin{\tikz[remember picture]{\node[text width=\marginparwidth,inner sep=0pt] (s86right) {\deleted{Sous cette forme, la condition implique la fidélité}};}}
% The cancelled right-margin note has a left-pointing arrow to the union condition.
i.e.
\[
\boxed{\forall x\in \tikz[remember picture,baseline=(s86xthree.base)]{\node[inner sep=0pt] (s86xthree) {$X$};}-|F|,\quad\exists Z\text{ avec }x\in|Z|^{\circ}\subset X\setminus|F|}
\]
\figure{\begin{tikzpicture}[remember picture,overlay]
\draw[->,gray] (s86note.east) to[out=15,in=-120] (s86xone.south west);
\draw[->] (s86note.east) to[out=25,in=-110] (s86xtwo.south west);
\draw[->] (s86note.east) to[out=0,in=90] (s86xthree.north);
\draw[->] (s86right.west) to[out=190,in=0] (s86cond.east);
\end{tikzpicture}}
\end{samepage}
Ceci posé, on trouve

\uline{Proposition} Soit $\unsure{\Phi|\mathscr P}\subset F$. Alors
\[
\varphi(\operatorname{supp}^{\circ}(\Phi))=|\Phi|^{\circ}
\left(\overset{\mathrm{def}}{=}\bigcup_{X\in\Phi}|X|^{\circ}\right)
\tag*{(*)}
\]

\uuline{Dém} Comme $\deleted{\unsure{X\in}}\Phi\subset\operatorname{supp}^{\circ}(\Phi)$, on a $\deleted{\varphi(\Phi)\subset}\varphi(\Phi)\overset{\mathrm{def}}{=}|\Phi|^{\circ}\subset\varphi(\operatorname{supp}^{\circ}(\Phi))$, i.e. on a $\supset$. Prouvons $\subset$, \deleted{Soit $Y\in\operatorname{supp}^{\circ}\Phi$, i.e.}
% slanted
\margin{expliciter d'abord la relation $\Phi\longrightarrow|\Phi|^{\circ}$, en termes de la figure image de $F$ dans $\unsure{P|\mathscr P}$\ldots}
\[
\forall Y\in\operatorname{supp}^{\circ}(\Phi)\Longrightarrow |Y|^{\circ}\subset|\Phi|^{\circ}\qquad\text{i.e.}
\]
\[
\begin{aligned}
|Y|^{\circ}\not\subset|\Phi|^{\circ}&\Longrightarrow Y\notin\operatorname{supp}^{\circ}(\Phi),\quad\text{i.e. }\exists Z\text{ avec}\\
\text{i.e. }|Y|^{\circ}\cap\complement|\Phi|^{\circ}\neq\varnothing
&\qquad Z\mathrel{|\circ|}\Phi,\quad\text{et }Z\mathrel{\overline{|\circ|}}Y,\quad\text{i.e.}\\
&|Z|^{\circ}\cap|\Phi|^{\circ}=\varnothing,\quad|Z|^{\circ}\cap|Y|^{\circ}\neq\varnothing.
\end{aligned}
\]
Distinguons deux cas

cas $|Y|^{\circ}\cap\bigl(\underbrace{|F|\setminus|\Phi|^{\circ}}_{=|F\setminus\Phi|^{\circ}}\bigr)\neq\varnothing$, \deleted{\unsure{Si $x$ est dans le}}
\[
\exists Z\in F\setminus\Phi,\quad\text{tel que }|Y|^{\circ}\cap|Z|^{\circ}\neq\varnothing.\quad\text{OK}
\]
cas $|Y|^{\circ}\cap(P\setminus|F|)\neq\varnothing$, alors on applique l'hypothèse de pleine fidélité à $F$.
\clearpage
\scanpage{87}
\mspage{78}
% archive sheet number: 87

% A vertical rule runs beside Corollaire 1.
\begin{samepage}
\noindent\tikz[remember picture,baseline=(s87top.base)]{\node[inner sep=0pt] (s87top) {\uline{Corollaire 1}};} L'application
\[
\unsure{\Phi|\mathscr P}\longmapsto\operatorname{supp}^{\circ}(\Phi):\quad\unsure{\mathscr P|\mathcal P}(F)\longrightarrow\overline\Sigma_M
\]
commute aux $\operatorname{Inf}$ et $\operatorname{Sup}$
\[
\begin{aligned}
\operatorname{supp}^{\circ}\left(\bigcap_i\Phi_i\right)&=\bigcap_i\operatorname{supp}^{\circ}(\Phi_i)\qquad I\neq\varnothing\\
\operatorname{supp}^{\circ}\left(\bigcup_i\Phi_i\right)&=\operatorname*{Sup}_i(\operatorname{supp}^{\circ}(\Phi_i))
\end{aligned}
\]
\margin{$\operatorname{sup}$ dans $\overline\Sigma_M$}
et de plus on a, si $\Phi\subset\Psi$
\[
\operatorname{supp}^{\circ}(\Psi\setminus\Phi)=\operatorname{supp}^{\circ}\Psi\cap\complement\operatorname{supp}^{\circ}\tikz[remember picture,baseline=(s87bottom.base)]{\node[inner sep=0pt] (s87bottom) {$\Phi$};}
\]
\figure{\begin{tikzpicture}[remember picture,overlay]
\path (s87top.west) ++(-0.5cm,0) coordinate (s87rule);
\draw (s87rule) -- (s87rule |- s87bottom.south);
\end{tikzpicture}}
\end{samepage}

\uline{Corollaire 2} \deleted{Dans} \inserted{Considérons} l'image $\Sigma_{M,F}$ de $\mathscr P(F)$ (ensemble des supports constructibles \uline{subordonnés} à $F$). Cette image est stable \inserted{dans $\overline\Sigma_M$} par $\operatorname{inf}$ (familles non vides) et $\operatorname{Sup}$, et complémentation relative i.e. opération $A\cap\complement B$. De plus, l'application $\deleted{A\unsure{-|\mapsto|=} |A|^{\circ}}$
\diagram{\begin{tikzcd}[ampersand replacement=\&,column sep=large,row sep=small]
A \arrow[r,mapsto] \& |A|^{\circ}=\displaystyle\bigcup_{X\in A}|X|^{\circ} \\
\Sigma_{M,F} \arrow[r] \arrow[d,phantom,"\cap" description] \& \mathscr P(|F|)\subset\mathscr P(P) \\
\overline\Sigma_M \&
\end{tikzcd}}
commute à ces opérations.

\uline{Terminologie} Un ens de \unsure{supp.\textsuperscript{ts}|supports} constructibles \inserted{(pour $\unsure{\mathscr F|\mathcal F}$)} est dit ``admissible'' \inserted{(pour $\mathscr F$)}, s'il existe $F\in\mathscr F$, \unsure{tels|tel} qu'\unsure{ils|il} soient tous \deleted{dans l'image de $\mathscr P(F)\unsure{-|\longrightarrow}\overline\Sigma_M$.} subordonnés à $F$.

\uline{Scholie} Pour les familles admissibles de supports constructibles, les opérations $\operatorname{Sup}.$ $\operatorname{Inf}$, complémentaire relatif \unsure{ont ttes|ont toutes}
\clearpage
\scanpage{88}
\mspage{79}
% archive sheet number: 88

les bonnes propriétés \deleted{\unsure{habit.|habit}} ensemblistes habituelles, notamment
\[
A\cap\operatorname*{Sup}_i B_i=\operatorname*{Sup}_i A\cap B_i
\]
\[
A\setminus(A\setminus B)=\deleted{\unsure{A}}B\qquad(\text{si }B\subset A)
\]
De plus \deleted{\unsure{A et|Soit}}
\[
\begin{aligned}
A\cap B=\varnothing&\Longleftrightarrow A\mathrel{|\circ|}B\ ;\\
&\text{(et pas seulement }\Leftarrow\text{)}
\end{aligned}
\]
Toute pathologie disparaît.

\uline{Question} Soit $M$ un magmien quelconque, sans supposer l'existence d'une fonction support externe pleinement fidèle. Ce scholie sur les familles admissibles de \unsure{fonctions|fonction} constructibles, \deleted{\unsure{et notamment}} \inserted{i.e.} l'énoncé-clef du cor.\ 1, sont-ils néanmoins valables ?
\margin{\uuline{Non}}

Aussi : \unsure{exemple|exple} de fonction support fidèle qui ne soit pleinement fidèle ??

\uline{Axiome} \uline{\deleted{\unsure{A|B}} supp 1} (``axiome des supports'').

\noindent\vrule\hspace{0.15em}\vrule\hspace{0.6em}\begin{minipage}{0.92\linewidth}
Soient : $F\in\unsure{\mathscr F|\mathcal F},\quad\unsure{\Phi|\mathscr P}\subset F,\quad\Phi'=F\setminus\Phi$.

Alors :
\[
\operatorname{supp}^{\circ}(F)\cap\deleted{\unsure{\complement_{\Phi}|\complement_{\Sigma}}}\operatorname{cosupp}^{\circ}\Phi=\operatorname{supp}^{\circ}(\Phi')
\]
en d'autres termes : $\forall Z\in M$
\[
\deleted{\unsure{\text{Si}}}\ Z\in\operatorname{supp}^{\circ}(F),\ \deleted{\text{et}}Z\mathrel{|\circ|}\Phi\ \Longrightarrow\ Z\in\operatorname{supp}^{\circ}(\Phi')
\]
\end{minipage}

\uline{NB} Cet axiome ne fait intervenir que $\leq$, $|\circ|$ sur $M$, et non pas $\overset{\circ}{\ll}$.
\clearpage
\scanpage{89}
\mspage{80}
% archive sheet number: 89

C'est le premier axiome de nature délicate, sur lequel nous \unsure{tombions|tombons}. Il est automatiquement satisfait pour les magmiens fidèles (car $\unsure{\operatorname{Ouv}|\operatorname{Oud}|\operatorname{Oub}}$ y est une fonction support ext. pleinement fidèle), plus \unsure{gén\textsuperscript{t}|génl\textsuperscript{t}} chaque fois qu'on a une fonction support externe pleinement fidèle. Le premier cas intéressant où il n'est pas sûr qu'on ait une telle fonction support, est celui du magmien associé à un ens ordonné, qu'il serait temps d'introduire ici en forme et d'examiner sous l'optique des supports.

Je voudrais donc à présent
\begin{enumerate}[label=\arabic*\textsuperscript{o}),leftmargin=*]
\item regarder le cas du magmien \deleted{\unsure{associé}} défini par un ens ordonné
\item Examiner le cas du magmien défini par $\leq$, $|\circ|$ seulement ($\overset{\circ}{\ll}$ étant défini via $\unsure{\leq|\ll}$, $|\circ|$, \deleted{\unsure{impliquant|indiquant} \unsure{une} \illegible{1}} \inserted{\unsure{entre des cinq.\ relations|entre diverses relations}} permettant de donner un sens à la relation $Y^{\circ}\subset X^{\circ}$ (en définissant $X^{\circ}=\unsure{\operatorname{supp}^{\circ}|\operatorname{Supp}^{\circ}}(X)\ldots$), donc de définir en conséquence $\overset{\circ}{\ll}$ \ldots
\item Voir si dans ces magmiens ``disjonctifs'', on n'aurait pas des exemples évidents où l'axiome des supports serait en défaut\ldots$^*$
\end{enumerate}
% slanted
\margin{$^*$ \uline{Oui} !}
